\section{Conclusion}\label{sec:conclusion}

Motivated by enterprise generative-AI services, we study sunk integration before future tariff architecture is chosen. The formal primitives are not AI-specific and can also represent other digital services with sunk integration and noncommitted future tariffs. On strict $\mathcal R^+$, anticipated metering produces a high-use installed state $h_M$ satisfying $h_0<h_M<h_F$, while the provider's gain from metering rises with the installed high-use mass. One fixed-state tariff threshold therefore becomes two self-consistency thresholds, $\mu_M<\mu_F$. Pure metering remains self-consistent at and below $\mu_M$, pure flat pricing at and above $\mu_F$, and only the strict interval between them requires the unique aggregate mixed resolution. If integration is architecture-insensitive, the two thresholds coincide.

The welfare result is narrower. Holding the installed base fixed, the provider meters over a larger activation-cost region than a social evaluator comparing the same decentralized allocations. With endogenous integration, the restricted architecture-welfare comparison is sign-indeterminate.

The candidate contribution is the specific reciprocal feedback from anticipated architecture to sunk integration and back to the relative profitability of that same architecture. It does not claim a global result outside strict $\mathcal R^+$, arbitrary-concave-demand generality, novelty for the generic expectation--participation--repricing sequence, or an AI-specific theorem. The Stage-11 prior-art repair also records a remaining source limitation for the complete published version of Hagiu (2006); journal positioning and whether to retain a GenAI-forward title are therefore deferred rather than treated as settled by the theory. Measuring the strength of the feedback in enterprise AI markets is left for empirical work.
